\documentclass[11pt]{amsart}
\usepackage[T1]{fontenc}
\usepackage{lmodern,microtype,amsmath,amssymb,amsthm,mathtools,enumitem}
\usepackage[margin=1.1in]{geometry}
\usepackage[hidelinks]{hyperref}
\usepackage{xurl}

\newtheorem{theorem}{Theorem}[section]
\newtheorem{lemma}[theorem]{Lemma}
\newtheorem{corollary}[theorem]{Corollary}
\newtheorem{proposition}[theorem]{Proposition}
\theoremstyle{definition}
\newtheorem{definition}[theorem]{Definition}
\theoremstyle{remark}
\newtheorem{remark}[theorem]{Remark}
\numberwithin{equation}{section}

\newcommand{\C}{\mathbb{C}}
\newcommand{\R}{\mathbb{R}}
\newcommand{\Hp}{\mathbb{C}^+}
\newcommand{\bp}{\boxplus}
\newcommand{\eps}{\varepsilon}
\DeclareMathOperator{\Log}{Log}

\hypersetup{pdftitle={Only point masses are freely strongly unimodal},pdfauthor={Yitzchak Shmalo}}

\title{Only point masses are freely strongly unimodal}
\author{Yitzchak Shmalo}
\address{Einstein Institute of Mathematics, The Hebrew University of Jerusalem, Jerusalem, Israel}
\email{yitzchak.shmalo@gmail.com}
\date{6 November 2026}
\subjclass[2020]{46L54, 60E05, 30C80}
\keywords{free additive convolution, unimodality, strong unimodality, subordination, Schwarz--Pick lemma}

\begin{document}
\begin{abstract}
A probability measure $\mu$ on $\R$ is called freely strongly unimodal if $\mu\bp\nu$ is unimodal for every unimodal probability measure $\nu$. Hasebe and Ueda showed that a measure with finite variance has this property only if it is a point mass, and asked whether any measure other than a point mass has it. We show that none does. For every probability measure $\mu$ that is not a point mass, the free convolution of $\mu$ with $\frac1{100}\delta_0+\frac{99}{100}\,U[0,L]$, where $U[0,L]$ is the uniform law on $[0,L]$, fails to be unimodal for all sufficiently large $L$. No moment condition is needed. The proof works with a small dilation of $\mu$, describes its inverse reciprocal Cauchy transform on a disc in the lower half-plane, and controls the resulting perturbation of the density by the Schwarz--Pick lemma.
\end{abstract}
\maketitle

\section{Introduction}\label{sec:intro}

A Borel probability measure $\nu$ on $\R$ is \emph{unimodal} if there is $a\in\R$ such that
\[
\nu=\nu(\{a\})\,\delta_a+f(x)\,dx,
\]
where $f$ is nondecreasing on $(-\infty,a)$ and nonincreasing on $(a,\infty)$. In classical probability, $\mu$ is strongly unimodal if $\mu*\nu$ is unimodal for every unimodal $\nu$; by Ibragimov's theorem these are exactly the measures that are absolutely continuous, supported on an interval, with a log-concave density (or point masses), so the normal laws are examples \cite[Theorem~52.3]{Sa13}. Hasebe and Ueda \cite{HU18} introduced the free analogue.

\begin{definition}[{\cite[Definition~3.8]{HU18}}]
A probability measure $\mu$ on $\R$ is \emph{freely strongly unimodal} if $\mu\bp\nu$ is unimodal for all unimodal distributions $\nu$ on $\R$.
\end{definition}

Here $\bp$ is the free additive convolution, defined for arbitrary Borel probability measures by Bercovici and Voiculescu \cite{BV93}. A point mass $\delta_a$ is freely strongly unimodal, since $\delta_a\bp\nu$ is the translate of $\nu$ by $a$. Hasebe and Ueda proved that semicircle laws are not freely strongly unimodal \cite[Lemma~3.9]{HU18}, and deduced from the free central limit theorem that no measure with finite variance other than a point mass is \cite[Theorem~3.10]{HU18}. They asked:

\begin{quote}
\emph{Problem 3.11 of \cite{HU18}.} Does there exist a freely strongly unimodal probability measure that is not a delta measure?
\end{quote}

The answer is no.

\begin{theorem}\label{thm:main}
Let $\mu$ be a Borel probability measure on $\R$ that is not a point mass, and for $L>0$ let
\[
\nu_L=\tfrac1{100}\,\delta_0+\tfrac{99}{100}\,U[0,L].
\]
Each $\nu_L$ is unimodal, and there is $L_0=L_0(\mu)$ such that $\mu\bp\nu_L$ is not unimodal for every $L>L_0$. Consequently the freely strongly unimodal probability measures on $\R$ are exactly the point masses.
\end{theorem}

The finite-variance argument of \cite{HU18} rescales $\mu^{\bp n}$ towards a semicircle law and so needs the second moment. Here we scale $\mu$ the other way. With $\eps=1/L$, the dilation $D_\eps$ (the push-forward under $x\mapsto\eps x$) carries $\mu\bp\nu_L$ to $D_\eps\mu\bp\kappa$, where
\begin{equation}\label{eq:kappa}
\kappa=p\,\delta_0+q\,U[0,1],\qquad p=\tfrac1{100},\quad q=\tfrac{99}{100}.
\end{equation}
The law $\kappa$ has an atom at its mode $0$ and the constant density $q$ on $(0,1)$. As $\eps\to0$ the measure $D_\eps\mu$ tends to $\delta_0$, and $D_\eps\mu\bp\kappa$ is a small perturbation of $\kappa$. The point is that the perturbation of the density on $(0,1)$ is not flat: to first order it is governed by a single holomorphic function, with values in the lower half-plane, evaluated along the curve $x\mapsto G_\kappa(x+i0)$, and the Schwarz--Pick lemma forces this function to have comparable size at two points of the curve. The curvature of the curve near $x=1/100$, coming from the atom, then produces a dip in the density that no tail behaviour of $\mu$ can cancel. Section~\ref{sec:prelim} collects the analytic facts we use, Section~\ref{sec:proof} proves Theorem~\ref{thm:main}, and Section~\ref{sec:numerics} reports a numerical check.

\section{Preliminaries}\label{sec:prelim}

For a Borel probability measure $\eta$ on $\R$ and $z$ in the upper half-plane $\Hp$ put
\[
G_\eta(z)=\int_\R\frac{\eta(ds)}{z-s},\qquad F_\eta(z)=\frac1{G_\eta(z)},\qquad h_\eta(z)=F_\eta(z)-z .
\]
The function $F_\eta$ maps $\Hp$ to $\Hp$.

\begin{lemma}\label{lem:F}
$\operatorname{Im}F_\eta(z)\ge\operatorname{Im}z$ for $z\in\Hp$, with strict inequality at every $z$ when $\eta$ is not a point mass.
\end{lemma}

\begin{proof}
We have $\operatorname{Im}G_\eta(z)=-\operatorname{Im}z\int|z-s|^{-2}\,\eta(ds)$, so
\[
\operatorname{Im}F_\eta(z)=\frac{-\operatorname{Im}G_\eta(z)}{|G_\eta(z)|^2}=\operatorname{Im}z\cdot\frac{\int|z-s|^{-2}\,\eta(ds)}{\bigl|\int(z-s)^{-1}\,\eta(ds)\bigr|^2}.
\]
By the Cauchy--Schwarz inequality the last quotient is at least $1$, with equality only if $s\mapsto(z-s)^{-1}$ is $\eta$-almost everywhere constant, that is, only if $\eta$ is a point mass.
\end{proof}

We use the analytic subordination theorem for free additive convolution of arbitrary Borel probability measures \cite{Bi98,BB07,Vo93}: if $\sigma=\lambda\bp\kappa$, there are analytic maps $\omega_\lambda,\omega_\kappa\colon\Hp\to\Hp$ with
\begin{equation}\label{eq:sub}
F_\lambda(\omega_\lambda(z))=F_\kappa(\omega_\kappa(z))=F_\sigma(z),\qquad \omega_\lambda(z)+\omega_\kappa(z)=z+F_\sigma(z).
\end{equation}
We need that \eqref{eq:sub} determines $F_\sigma(z)$ pointwise.

\begin{lemma}\label{lem:unique}
Fix $z\in\Hp$. If $\omega_1,\omega_2,u\in\Hp$ satisfy $F_\lambda(\omega_1)=F_\kappa(\omega_2)=u$ and $\omega_1+\omega_2=z+u$, then $u=F_\sigma(z)$.
\end{lemma}

\begin{proof}
From the two equations, $\omega_1=z+h_\kappa(\omega_2)$ and $\omega_2=z+h_\lambda(\omega_1)$, so $\omega_2$ is a fixed point of $T_z(v)=z+h_\lambda(z+h_\kappa(v))$. By Lemma~\ref{lem:F}, $h_\kappa$ and $h_\lambda$ map $\Hp$ into $\{\operatorname{Im}\ge0\}$, so $T_z$ is a holomorphic map of $\Hp$ into $\{\operatorname{Im}v\ge\operatorname{Im}z\}$; in particular it is not the identity. A holomorphic self-map of $\Hp$ other than the identity has at most one fixed point in $\Hp$, by the Schwarz lemma transported to the half-plane. The subordination functions provide one fixed point, $\omega_\kappa(z)$, so $\omega_2=\omega_\kappa(z)$ and $u=F_\kappa(\omega_2)=F_\sigma(z)$.
\end{proof}

\begin{lemma}\label{lem:inverse}
Let $\mu$ be a probability measure that is not a point mass, $\lambda_\eps=D_\eps\mu$, and let $\mathcal D$ be an open disc whose closure lies in the lower half-plane. Put $U=\{1/w:w\in\overline{\mathcal D}\}$, a compact subset of $\Hp$. There are a bounded open set $V$ with $U\subset V\subset\overline V\subset\Hp$ and $\eps_0>0$ such that for $0<\eps<\eps_0$:
\begin{enumerate}[label=\textup{(\roman*)}]
\item $F_{\lambda_\eps}$ is injective on $V$, and there is a holomorphic $H_\eps$ on a neighbourhood of $U$ with values in $V$ and $F_{\lambda_\eps}\circ H_\eps=\mathrm{id}$;
\item the function $R_\eps(w)=H_\eps(1/w)-1/w$ is holomorphic on a neighbourhood of $\overline{\mathcal D}$, satisfies $\operatorname{Im}R_\eps<0$ there, and $\sup_{\overline{\mathcal D}}|R_\eps|\to0$ as $\eps\to0$.
\end{enumerate}
\end{lemma}

\begin{proof}
For $z\in\Hp$, $G_{\lambda_\eps}(z)=\int(z-\eps s)^{-1}\mu(ds)$, and $|z-\eps s|\ge\operatorname{Im}z$, so by dominated convergence $G_{\lambda_\eps}(z)\to1/z$ pointwise. The family is bounded by $1/\operatorname{Im}z$, so by Vitali's theorem the convergence is locally uniform on $\Hp$, together with all derivatives, and the same holds for $F_{\lambda_\eps}\to\mathrm{id}$.

Let $\delta>0$ be less than half the distance from $U$ to $\R$, and let $V$ be the open $2\delta$-neighbourhood of the convex hull of $U$; it is convex, bounded, and $\overline V\subset\Hp$. Choose $\eps_0$ so that $|F_{\lambda_\eps}'-1|<1/2$ and $|F_{\lambda_\eps}-\mathrm{id}|<\delta/2$ on $\overline V$ for $\eps<\eps_0$. Integrating along segments gives $\operatorname{Re}\frac{F_{\lambda_\eps}(\omega)-F_{\lambda_\eps}(\omega')}{\omega-\omega'}>\frac12$ on $V$, so $F_{\lambda_\eps}$ is injective there. For $u'$ within $\delta/2$ of $U$, Rouch\'e's theorem on the circle $|\omega-u'|=\delta$ (which lies in $V$) shows that $F_{\lambda_\eps}(\omega)=u'$ has exactly one solution in that disc; this defines $H_\eps$, which is holomorphic as the inverse of an injective holomorphic map. Since $H_\eps(u')-u'=H_\eps(u')-F_{\lambda_\eps}(H_\eps(u'))$, we get $|H_\eps(u')-u'|\le\sup_{V}|F_{\lambda_\eps}-\mathrm{id}|\to0$.

This gives (i) and the last assertion of (ii). For the sign, let $u=1/w$ with $w$ near $\overline{\mathcal D}$ and $\omega=H_\eps(u)\in\Hp$. Since $\lambda_\eps$ is not a point mass, Lemma~\ref{lem:F} gives $\operatorname{Im}u=\operatorname{Im}F_{\lambda_\eps}(\omega)>\operatorname{Im}\omega$, so $\operatorname{Im}R_\eps(w)=\operatorname{Im}(\omega-u)<0$.
\end{proof}

No moment of $\mu$ enters Lemma~\ref{lem:inverse}, and no global inverse of $F_{\lambda_\eps}$ is used.

\section{Proof of Theorem~\ref{thm:main}}\label{sec:proof}

\subsection{The reference measure and the disc}
Let $\kappa$ be as in \eqref{eq:kappa}. Then
\[
G_\kappa(z)=\frac pz+q\bigl(\log z-\log(z-1)\bigr),
\]
and for $0<x<1$ its boundary value from $\Hp$ is $G_\kappa(x+i0)=g(x)-ic$, where
\[
g(x)=\frac px+q\log\frac{x}{1-x},\qquad c=\pi q .
\]
For $z$ near a point of $(0,1)$ put
\[
A(z)=\frac pz+q\bigl(\Log z-\Log(1-z)-i\pi\bigr)
\]
with the principal branch of $\Log$. Then $A$ is holomorphic on a neighbourhood of $(0,1)$, agrees with $G_\kappa$ on its intersection with $\Hp$, and $A(x)=g(x)-ic$ for real $x\in(0,1)$. In particular $A'(x)=g'(x)$ is real there. Direct differentiation gives
\begin{equation}\label{eq:g}
g'(x)=\frac{x-p}{x^2(1-x)},\qquad g''(x)=\frac{2p}{x^3}+q\Bigl(\frac1{(1-x)^2}-\frac1{x^2}\Bigr).
\end{equation}

Let $x_1=p=1/100$ and $x_2=1/4$. By \eqref{eq:g},
\begin{equation}\label{eq:gvals}
g'(x_1)=0,\qquad g''(x_1)=\frac{1000000}{99},\qquad g'(x_2)=\frac{128}{25},\qquad g''(x_2)=-\frac{64}{5}.
\end{equation}
Put $w_j=g(x_j)-ic$ and
\[
b=g(x_2)-10i,\qquad r=\tfrac{99}{10},\qquad \mathcal D=\{w:|w-b|<r\}.
\]
The closure of $\mathcal D$ lies in $\{\operatorname{Im}w\le-1/10\}$. Since $3<\pi<4$ we have $2.97<c<3.96$. Let $d=g(x_1)-g(x_2)=\frac{24}{25}-\frac{99}{100}\log33$; since $1<\log 33<4$, $|d|<3$. Now $w_2-b=i(10-c)$ and $w_1-b=d+i(10-c)$, so
\[
|w_2-b|<7.03<r,\qquad |w_1-b|^2<9+7.03^2<r^2,
\]
and both $w_j$ lie in $\mathcal D$. Write $\Delta=r^2-(10-c)^2$. Then
\[
\Delta>98.01-7.03^2=48.5891,\qquad \frac{2r}{\Delta}<\frac{19.8}{48.5891}<\frac5{12}.
\]
The pseudohyperbolic distance of $w_1$ and $w_2$ in $\mathcal D$ is
\[
\rho=\frac{r|w_1-w_2|}{\bigl|r^2-(w_1-b)\overline{(w_2-b)}\bigr|}=\frac{r|d|}{|\Delta+id(10-c)|}\le\frac{r|d|}{\Delta}<\frac{29.7}{48.5891}<\frac23 .
\]
(Numerically $\Delta\approx50.54$, $d\approx-2.5015$ and $\rho\approx0.464$.)

\subsection{The perturbation and its size}
Fix $\mu$, not a point mass, and apply Lemma~\ref{lem:inverse} to $\lambda_\eps=D_\eps\mu$ and the disc $\mathcal D$. Put
\[
a_\eps=\operatorname{Re}R_\eps(w_2),\qquad t_\eps=-\operatorname{Im}R_\eps(w_2)>0,\qquad B_\eps=R_\eps-a_\eps .
\]
Then $t_\eps\to0$, $B_\eps$ maps $\mathcal D$ into the lower half-plane, and $B_\eps(w_2)=-it_\eps$. Write $t=t_\eps$, $B=B_\eps$ and $K=-\operatorname{Im}B>0$ on $\mathcal D$.

\begin{lemma}\label{lem:SP}
For every compact $M\subset\mathcal D$ and every $k\ge0$ there is $C_{M,k}$, independent of $\eps$, with $|B^{(k)}|\le C_{M,k}\,t$ on $M$. Moreover
\[
K(w_1)>\tfrac15\,t,\qquad K(w_2)=t,\qquad |\partial_xK(w_2)|<\tfrac5{12}\,t,
\]
where $\partial_x$ is the derivative in the real direction.
\end{lemma}

\begin{proof}
Since $\operatorname{Im}B<0$ and $t>0$, the function $\phi=(B+it)/(B-it)$ maps $\mathcal D$ into the unit disc and $\phi(w_2)=0$. By the Schwarz--Pick lemma, $|\phi(w)|\le\rho_{\mathcal D}(w,w_2)$ and $|\phi'(w_2)|\le r/(r^2-|w_2-b|^2)=r/\Delta$. Solving for $B$,
\[
B=-it\,\frac{1+\phi}{1-\phi},\qquad K=t\,\frac{1-|\phi|^2}{|1-\phi|^2}\ \ge\ t\,\frac{1-|\phi|}{1+|\phi|}.
\]
On a compact $M\subset\mathcal D$ we have $\sup_M|\phi|\le\sup_{w\in M}\rho_{\mathcal D}(w,w_2)<1$, which bounds $B$ by a multiple of $t$ on $M$; Cauchy's estimates on a slightly larger compact set bound the derivatives. At $w_1$, $K(w_1)\ge t(1-\rho)/(1+\rho)>t\cdot\frac{1/3}{5/3}=t/5$. Finally $B'(w_2)=-2it\,\phi'(w_2)$, so $|\partial_xK(w_2)|=|\operatorname{Im}B'(w_2)|\le|B'(w_2)|\le 2rt/\Delta<\frac5{12}t$.
\end{proof}

\subsection{The density of the convolution}
Let $\tilde\lambda_\eps$ be the translate of $\lambda_\eps$ by $-a_\eps$, so that $F_{\tilde\lambda_\eps}(\omega)=F_{\lambda_\eps}(\omega+a_\eps)$, and let $\sigma_\eps=\tilde\lambda_\eps\bp\kappa$. Fix $j\in\{1,2\}$ and a closed disc $N_j$ centred at $x_j$, small enough that $A$ is holomorphic on a neighbourhood of $N_j$ and $A(N_j)$ is contained in a compact set $M\subset\mathcal D$ with $\operatorname{dist}(A(N_j),\partial M)>0$.

For $z\in N_j$ consider the equation
\begin{equation}\label{eq:W}
W=A\bigl(z-B(W)\bigr).
\end{equation}
By Lemma~\ref{lem:SP}, $B$ and $B'$ are $O(t)$ on $M$, so for small $\eps$ the right side of \eqref{eq:W} maps the closed disc of radius $C t$ about $A(z)$ into itself and is a contraction there, for a suitable $C$ independent of $\eps$ and $z$, after shrinking $N_j$ slightly. The fixed point $W_\eps(z)$ is holomorphic in $z$ (the iterates converge uniformly), lies in $M$, and satisfies $W_\eps=A+O(t)$ uniformly on $N_j$.

\begin{lemma}\label{lem:ident}
$W_\eps(z)=G_{\sigma_\eps}(z)$ for $z\in N_j\cap\Hp$.
\end{lemma}

\begin{proof}
Let $z\in N_j\cap\Hp$, $W=W_\eps(z)$ and $u=1/W\in U$. Put
\[
\omega_1=H_\eps(u)-a_\eps=u+B(W),\qquad \omega_2=z-B(W).
\]
Then $\omega_1\in\Hp$ because $H_\eps(u)\in V$, and $\operatorname{Im}\omega_2=\operatorname{Im}z+K(W)>0$. By Lemma~\ref{lem:inverse}, $F_{\tilde\lambda_\eps}(\omega_1)=F_{\lambda_\eps}(H_\eps(u))=u$. Since $\omega_2\in\Hp$ is near $x_j$, $A(\omega_2)=G_\kappa(\omega_2)$, so \eqref{eq:W} says $F_\kappa(\omega_2)=1/W=u$. Finally $\omega_1+\omega_2=z+u$. Lemma~\ref{lem:unique} gives $u=F_{\sigma_\eps}(z)$, that is $W=G_{\sigma_\eps}(z)$.
\end{proof}

Thus $G_{\sigma_\eps}$ extends holomorphically across the real interval $N_j\cap\R$, and by Stieltjes inversion $\sigma_\eps$ is absolutely continuous there with the real-analytic density
\[
f_\eps(x)=-\frac1\pi\operatorname{Im}W_\eps(x),\qquad x\in N_j\cap\R .
\]

\subsection{First-order expansion}
Expanding \eqref{eq:W} with $W_\eps=A+O(t)$ and $B,B'=O(t)$ on $M$,
\[
W_\eps(z)=A(z)-A'(z)\,B(W_\eps(z))+O(t^2)=A(z)-A'(z)\,B(A(z))+O(t^2)
\]
uniformly on $N_j$, with constants independent of $\eps$. The remainder is holomorphic, so by Cauchy's estimates its derivative is also $O(t^2)$ on a smaller concentric disc. For real $x$ we have $A(x)=g(x)-ic$ and $A'(x)=g'(x)\in\R$, hence
\[
f_\eps(x)=q-\frac1\pi\,g'(x)\,K\bigl(A(x)\bigr)+O(t^2).
\]
As $x$ moves along the real axis, $A(x)$ moves horizontally with speed $g'(x)$, so $\frac{d}{dx}K(A(x))=g'(x)\,\partial_xK(A(x))$ and
\begin{equation}\label{eq:fprime}
f_\eps'(x)=-\frac1\pi\Bigl(g''(x)\,K\bigl(A(x)\bigr)+g'(x)^2\,\partial_xK\bigl(A(x)\bigr)\Bigr)+O(t^2).
\end{equation}

\subsection{Conclusion}
At $x_1$, \eqref{eq:gvals}, \eqref{eq:fprime} and Lemma~\ref{lem:SP} give
\[
f_\eps'(x_1)=-\frac{g''(x_1)}{\pi}K(w_1)+O(t^2)\le-\frac{200000}{99\,\pi}\,t+O(t^2).
\]
At $x_2$,
\[
f_\eps'(x_2)\ge\frac1\pi\Bigl(\frac{64}5-\Bigl(\frac{128}{25}\Bigr)^2\frac5{12}\Bigr)t-O(t^2)=\frac{704}{375\,\pi}\,t-O(t^2).
\]
Since $t=t_\eps>0$ and $t_\eps\to0$, there is $\eps_1>0$ such that for $0<\eps<\eps_1$ the density $f_\eps$ of $\sigma_\eps$ is strictly decreasing near $x_1$ and strictly increasing near $x_2$. Then $\sigma_\eps$ is not unimodal: if $a$ were a mode, the density would be nondecreasing almost everywhere to the left of $a$, which forces $a\le x_1$, and nonincreasing to the right of $a$, which fails near $x_2>x_1$.

Finally, $D_\eps(\mu\bp\nu_{1/\eps})=D_\eps\mu\bp D_\eps\nu_{1/\eps}=\lambda_\eps\bp\kappa$, and $\sigma_\eps$ is the translate of this measure by $-a_\eps$. Translations and positive dilations preserve unimodality, so $\mu\bp\nu_L$ is not unimodal for $L>1/\eps_1$. Each $\nu_L$ is unimodal with mode $0$, so $\mu$ is not freely strongly unimodal. Together with the remark on point masses in the introduction this proves Theorem~\ref{thm:main}. \qed

\begin{remark}
The atom of $\kappa$ is what makes $g'$ vanish at an interior point of the support, with $g''$ large there. The proof gives no explicit $L_0(\mu)$, since $t_\eps$ is not quantified in terms of $\mu$.
\end{remark}

\section{A numerical check}\label{sec:numerics}

As an independent check of the sign pattern we computed the density of $\sigma\bp\kappa$ directly for small $\sigma$ of two kinds. For the semicircle law of variance $v$ the subordination equation reduces to $G(z)=G_\kappa\bigl(z-vG(z)\bigr)$, solved by fixed-point iteration at $z=x+10^{-13}i$; for the symmetric Bernoulli law on $\{\pm a\}$ we iterated $\omega\mapsto z+h(z+h_\kappa(\omega))$ with $h(w)=-a^2/w$. Central differences with step $10^{-5}$ gave the following.

\begin{center}
\begin{tabular}{lccc}
law of $\sigma$ & $f'(0.01)$ & $f'(0.25)$ & interior local minimum\\\hline
semicircle, variance $10^{-6}$ & $-9.99\times10^{-3}$ & $1.27\times10^{-5}$ & $x\approx0.020$\\
semicircle, variance $10^{-4}$ & $-8.66\times10^{-1}$ & $1.26\times10^{-3}$ & $x\approx0.020$\\
Bernoulli $\pm10^{-3}$ & $-9.99\times10^{-3}$ & $1.27\times10^{-5}$ & $x\approx0.020$
\end{tabular}
\end{center}

\noindent In each case the density has an interior local minimum just to the right of $x_1$, as the proof predicts. These measures have finite variance, so they are covered by \cite[Theorem~3.10]{HU18}; the computation only confirms the mechanism. The arithmetic of Section~\ref{sec:proof} (the values \eqref{eq:gvals}, the bounds on $\Delta$, $\rho$ and $2r/\Delta$, and the two slopes) was also checked in exact rational and 50-digit arithmetic.

\section*{Acknowledgements and funding}
The research presented in this paper was supported by the European Research
Council (ERC) under the European Union's Horizon Europe research and
innovation programme (grant agreement No.~101041711), by the Simons
Foundation as part of the Collaboration on the Mathematical and Scientific
Foundations of Deep Learning, by Heights Labs, by Convex Nexus Capital, by
the Israel Science Foundation (grant number 2258/19), and by the Israel
Science Foundation (ISF Grant 4101/25). Computational resources were
provided in part by Google Cloud through its education credits program
(credit 530550103).

\section*{Declaration of competing interest}
The author holds equity in Heights Labs and in Convex Nexus Capital and
has a management role at Convex Nexus Capital. These entities are included
in the funding declaration above.

\section*{Declaration of generative AI and AI-assisted technologies}
The paper comes out of Math Brain, an ongoing project on the use of AI systems for research mathematics. In October 2026 a GPT-family model (OpenAI), running in the Codex command-line tool, found the proof of Theorem~\ref{thm:main}; a Claude model (Anthropic) reviewed it. In November 2026 Claude Opus~5.5 (Anthropic), working in Claude Code, checked the proof and its constants, wrote and ran the computations of Section~\ref{sec:numerics}, and wrote the paper. The author directed the work and is responsible for the content of the paper and for decisions concerning its dissemination.

\section*{Data availability}
No data were used. The code for the computations of Section~\ref{sec:numerics} and for the check of the constants, with their outputs, is available at\\ \url{https://github.com/yspennstate/free-strong-unimodality-point-masses}.

\begin{thebibliography}{9}

\bibitem{BB07}
S.~T.~Belinschi and H.~Bercovici, A new approach to subordination results in free probability, \emph{J. Anal. Math.} \textbf{101} (2007), 357--365.

\bibitem{BV93}
H.~Bercovici and D.~V.~Voiculescu, Free convolution of measures with unbounded support, \emph{Indiana Univ. Math. J.} \textbf{42} (1993), 733--773.

\bibitem{Bi98}
P.~Biane, Processes with free increments, \emph{Math. Z.} \textbf{227} (1998), 143--174.

\bibitem{HU18}
T.~Hasebe and Y.~Ueda, Large time unimodality for classical and free Brownian motions with initial distributions, \emph{ALEA Lat. Am. J. Probab. Math. Stat.} \textbf{15} (2018), 353--374; arXiv:1710.08240.

\bibitem{Sa13}
K.~Sato, \emph{L\'evy Processes and Infinitely Divisible Distributions}, corrected paperback edition, Cambridge University Press, Cambridge, 2013.

\bibitem{Vo93}
D.~Voiculescu, The analogues of entropy and of Fisher's information measure in free probability theory, I, \emph{Comm. Math. Phys.} \textbf{155} (1993), 71--92.

\end{thebibliography}
\end{document}
